\section{RÚBRICA}

\begin{center}\renewcommand{\arraystretch}{1.5}\begin{tabular}{|l|l|c|}\hline\rowcolor[HTML]{CCCCCC}\multicolumn{1}{|c|}{\cellcolor[HTML]{CCCCCC}Tarea}                                                                                & \multicolumn{1}{c|}{\cellcolor[HTML]{CCCCCC}Puntuación} & Máximo \\ \hline
\textit{Presentación del informe}                                                                                                  &                                                         & 20\%   \\ \hline
\textit{Comprender las regiones de operación de los tiristores}                                                                    &                                                         & 30\%   \\ \hline
\textit{\begin{tabular}[c]{@{}l@{}}Identificación de las especificaciones de los tiristores desde\\ la hoja de datos\end{tabular}} &                                                         & 30\%   \\ \hline
\textit{Defensa de las conclusiones}                                                                                               &                                                         & 20\%   \\ \hline
\textbf{Total}                                                                                                                     &                                                         & 100\%  \\ \hline
\end{tabular}\end{center}


\newpage

\section{INTRODUCCIÓN}
\sangria{} El objetivo de este informe es relevar los parámetros de fabricación de los tiristores: SCR - DIAC - TRIAC. Se analizarán los mismos por medio del simulador LTSpite y por una serie de circuitos en protoboard para identificar las regiones de trabajo de estos dispositivos.

\subsection{Condiciones de trabajo y ensayo}

